\documentclass[10pt,letterpaper]{article}

\usepackage[utf8]{inputenc}
\usepackage[T1]{fontenc}
\usepackage[margin=0.78in]{geometry}
\usepackage{hyperref}
\usepackage{xcolor}
\usepackage{enumitem}
\usepackage{multicol}
\usepackage{titlesec}
\usepackage{lmodern}
\usepackage{float}
\usepackage{caption}
\usepackage{tikz}

\usetikzlibrary{calc,positioning}

\setlength{\parindent}{0pt}
\setlength{\parskip}{0.24em}
\setlength{\columnsep}{0.2in}
\renewcommand{\baselinestretch}{0.965}
\titlespacing*{\section}{0pt}{0.55em}{0.28em}
\titleformat{\section}{\normalsize\bfseries}{}{0em}{}

\definecolor{graphBlue}{HTML}{7D93B2}
\definecolor{llmOrange}{HTML}{D7A36A}
\definecolor{aiGreen}{HTML}{8FAA83}
\definecolor{ink}{HTML}{1F2937}
\definecolor{muted}{HTML}{6B7280}
\definecolor{accentGreen}{HTML}{0B8A4A}
\definecolor{refBlue}{HTML}{1F77B4}

\hypersetup{
    colorlinks=true,
    linkcolor=refBlue,
    urlcolor=refBlue,
    citecolor=refBlue
}

\newcommand{\urlScholar}{https://scholar.google.com/citations?user=Pb_UoYwAAAAJ}
\newcommand{\urlNSHE}{https://www.ijcai.org/proceedings/2020/190}
\newcommand{\urlHGSL}{https://doi.org/10.1609/aaai.v35i5.16600}
\newcommand{\urlHousE}{https://proceedings.mlr.press/v162/li22ab.html}
\newcommand{\urlGophormer}{https://arxiv.org/abs/2110.13094}
\newcommand{\urlMVSE}{https://doi.org/10.1007/978-3-030-86520-7_20}
\newcommand{\urlGSR}{https://doi.org/10.1145/3539597.3570455}
\newcommand{\urlGLEM}{https://openreview.net/forum?id=q0nmYciuuZN}
\newcommand{\urlGraphText}{https://arxiv.org/abs/2310.01089}
\newcommand{\urlGraphAny}{https://openreview.net/forum?id=1Qpt43cqhg}
\newcommand{\urlHAX}{https://arxiv.org/abs/2507.00449}
\newcommand{\urlPandemicLLM}{https://www.nature.com/articles/s43588-025-00798-6}
\newcommand{\urlRaftPPI}{https://openreview.net/forum?id=Dp1RM3gPg8}
\newcommand{\urlGeneZip}{https://arxiv.org/abs/2602.17739}

\newcommand{\topictag}[2]{%
  {\textcolor{#1!70!black}{\textbf{\scriptsize #2}}}%
}

\newcommand{\paperrow}[4]{%
  \noindent\href{#4}{\textcolor{refBlue}{\textbf{[#1]}}} \href{#4}{\textcolor{ink}{#2}} {\color{muted}|}~#3\par
}

\newcommand{\pubgroup}[1]{%
  \vspace{0.12em}{\color{ink}\textbf{#1}}\par\vspace{0.06em}%
}

\newcommand{\papermark}[3]{%
  \node[
    font=\tiny\bfseries,
    text=refBlue,
    fill=white,
    fill opacity=0.92,
    text opacity=1,
    inner xsep=0.7pt,
    inner ysep=0.1pt,
    rounded corners=1pt
  ] at #1 {\href{#2}{\textcolor{refBlue}{[#3]}}};
}

\begin{document}
\pagestyle{empty}
\small

\begin{center}
{\Large \textbf{Statement of Research}}\\
Jianan Zhao\hspace{1.6em}{\small\color{ink}\textbf{PhD Candidate at Mila-UdeM (Expected graduation: Jan 2027)}}
\end{center}

\section*{1. Research Expertise}
My research lies at the intersection of graph machine learning, foundation models, and AI for science. Across these areas, I have tried to build methods that are simple in formulation, scalable in practice, and faithful to the structure of complex data. A consistent theme in my work is that many real-world domains --- including scientific data, biological systems, and knowledge bases --- are not naturally iid tables or free text, but richly structured objects with relations, hierarchies, and geometry. I am therefore especially interested in models that can preserve this structure while still benefiting from the transfer, scale, and flexibility of modern foundation models. This perspective aligns naturally with Zitnik Lab's emphasis on geometry, knowledge-grounded learning, multimodal AI, and impactful applications in science and medicine. Figure~1 summarizes selected publications across these themes.

\begin{figure}[H]
\centering
\makebox[\linewidth][c]{%
\begin{minipage}[t]{0.35\linewidth}
\centering
\vspace{0pt}
\begin{tikzpicture}[scale=0.64]
  \def\r{2.18}
  \def\ra{1.45}
  \coordinate (G) at (-1.20,-0.08);
  \coordinate (L) at ( 1.50,1.04);
  \coordinate (A) at ( 1.42,-1.58);

  \fill[graphBlue!28, opacity=0.48] (G) circle (\r);
  \fill[llmOrange!34, opacity=0.42] (L) circle (\r);
  \fill[aiGreen!30, opacity=0.44] (A) circle (\ra);

  \draw[line width=0.9pt, graphBlue!75!black] (G) circle (\r);
  \draw[line width=0.9pt, llmOrange!78!black] (L) circle (\r);
  \draw[line width=0.9pt, aiGreen!72!black] (A) circle (\ra);

  \node[font=\bfseries\normalsize, text=graphBlue!65!black] at (-1.82,0.5) {Graph};
  \node[font=\bfseries\normalsize, text=llmOrange!72!black, align=center] at (1.78,1.82) {LLM / \\GFM};
  \node[font=\bfseries\normalsize, text=aiGreen!65!black] at (1.46,-1.80) {AI4S};

  \papermark{(-2.28,-0.18)}{\urlGSR}{8}
  \papermark{(-1.62,-0.18)}{\urlHousE}{9}
  \papermark{(-0.92,-0.18)}{\urlGophormer}{10}
  \papermark{(-2.12,-0.80)}{\urlMVSE}{11}
  \papermark{(-1.40,-0.80)}{\urlHGSL}{12}
  \papermark{(-0.68,-0.80)}{\urlNSHE}{13}

  \papermark{(-0.1,1.1)}{\urlGraphAny}{5}
  \papermark{(0.1,0.6)}{\urlGraphText}{6}
  \papermark{(0.3,0.1)}{\urlGLEM}{7}

  \papermark{(1.90,0.46)}{\urlHAX}{3}

  \papermark{(1.12,-0.68)}{\urlGeneZip}{1}
  \papermark{(1.67,-0.68)}{\urlRaftPPI}{2}
  \papermark{(2.22,-0.68)}{\urlPandemicLLM}{4}
\end{tikzpicture}
\end{minipage}\hspace{0.018\linewidth}%
\begin{minipage}[t]{0.57\linewidth}
\vspace{0pt}
\raggedright
\scriptsize
\setlength{\parskip}{0.03em}
\pubgroup{PhD Publications}
\paperrow{1}{arXiv26 \textit{GeneZip}}{\topictag{llmOrange}{LLM}/\topictag{aiGreen}{AI4S}}{\urlGeneZip}
\paperrow{2}{ICLR26 \textit{RaftPPI}}{\topictag{llmOrange}{LLM}/\topictag{aiGreen}{AI4S}}{\urlRaftPPI}
\paperrow{3}{NeurIPS25 \textit{HAX}}{\topictag{llmOrange}{LLM}}{\urlHAX}
\paperrow{4}{Nature Computational Science 2025 \textit{PandemicLLM}}{\topictag{llmOrange}{LLM}/\topictag{aiGreen}{AI4S}}{\urlPandemicLLM}
\paperrow{5}{ICLR25 \textit{GraphAny}}{\topictag{graphBlue}{Graph}/\topictag{llmOrange}{GFM}}{\urlGraphAny}
\paperrow{6}{arXiv23 \textit{GraphText}}{\topictag{graphBlue}{Graph}/\topictag{llmOrange}{LLM}}{\urlGraphText}
\paperrow{7}{ICLR23 \textit{GLEM}}{\topictag{graphBlue}{Graph}/\topictag{llmOrange}{LLM}}{\urlGLEM}
\pubgroup{Previous Publications}
\paperrow{8}{WSDM23 \textit{GSR}}{\topictag{graphBlue}{Graph}}{\urlGSR}
\paperrow{9}{ICML22 \textit{HousE}}{\topictag{graphBlue}{Graph}}{\urlHousE}
\paperrow{10}{arXiv21 \textit{Gophormer}}{\topictag{graphBlue}{Graph}}{\urlGophormer}
\paperrow{11}{ECML21 \textit{MVSE}}{\topictag{graphBlue}{Graph}}{\urlMVSE}
\paperrow{12}{AAAI21 \textit{HGSL}}{\topictag{graphBlue}{Graph}}{\urlHGSL}
\paperrow{13}{IJCAI20 \textit{NSHE}}{\topictag{graphBlue}{Graph}}{\urlNSHE}
\end{minipage}
}
\caption{\small Selected publications across Graph, LLM / GFM, and AI4S (only first-, second-, or equal-contribution author works). Full list: \href{\urlScholar}{Google Scholar}. Numbered markers and paper titles are clickable in the PDF.}
\end{figure}
\vspace{-0.8em}

My early work focused on structured and heterogeneous graphs. In \textit{Network Schema Preserving Heterogeneous Information Network Embedding} and \textit{Heterogeneous Graph Structure Learning}~\cite{nshe,hgsl}, I studied how to represent graph data when node and edge types carry different semantics, and how to improve graph learning by revising the input structure rather than treating the graph as fixed. These projects taught me that structure is not just a storage format but a source of inductive bias, and that useful models should expose abstractions that transfer across settings. I later extended this direction to graph transformers in \textit{Gophormer}~\cite{gophormer}, graph structure refinement in \textit{GSR}~\cite{gsr}, and knowledge-graph representation learning in \textit{HousE}~\cite{house}. Together, this work gave me a strong base in graph representation learning, scalable GNNs, and knowledge-aware machine learning.

During my PhD, I became increasingly interested in bridging graph learning with language models. In \textit{GLEM}~\cite{glem}, I developed a variational EM framework that couples a language model and a GNN on large text-attributed graphs without the cost of end-to-end joint training. The method showed that language and graph structure can improve one another through iterative co-training and achieved top Open Graph Benchmark performance. More importantly, \textit{GLEM} convinced me that graph-model/LLM interaction is a first-class modeling problem rather than a feature concatenation trick.

I pushed this interface further in \textit{GraphText}~\cite{graphtext}, asking whether graph problems could be translated into a language space that LLMs can operate on directly. By converting local graph structure into a graph-syntax tree and textual sequence, the framework enabled training-free, interactive graph reasoning and showed that the right graph-to-text interface can let an LLM outperform supervised GNN baselines on node classification without graph-specific training. In parallel, \textit{GraphAny}~\cite{graphany} asked the complementary question of fully-inductive generalization, where test graphs may have unseen structure, features, and label spaces. Together, these projects shaped my view that graph ML is moving toward reusable, structure-aware priors and interfaces rather than isolated task-specific models.

More recently, I have focused on foundation-model ideas for scientific problems where structured reasoning and systems constraints both matter. In \textit{PandemicLLM}~\cite{pandemicllm}, we reframed infectious disease forecasting as multimodal reasoning over text, policy, genomic surveillance, and epidemiological signals. In \textit{RaftPPI}~\cite{raftppi}, we developed a residue-level factorization strategy for proteome-scale protein interaction retrieval, turning expensive screening into retrieval-friendly computation. In \textit{GeneZip}~\cite{genezip}, we studied long-context DNA modeling through region-aware compression, using biological priors to allocate model capacity where it matters most. These projects broadened my research beyond graph benchmarks and reinforced my belief that foundation models are most meaningful when they become reliable tools for scientific reasoning and decision support.

Taken together, my prior work has prepared me to contribute to Zitnik Lab in several ways. I can bring experience in graph learning, graph-language interfaces, large-scale model design, and AI-for-science applications. Just as importantly, I bring a research style that tries to connect method development with a broader conceptual question: what should the interface between structure, language, and scientific knowledge look like? At the same time, I see Zitnik Lab as an environment where I can grow substantially. The lab's strengths in biomedical knowledge graphs, geometric learning, multimodal models, and scientific deployment are precisely the directions in which I want to push my research further. I view myself as a good match not because my current work already solves these problems, but because it gives me a strong starting point for contributing to them in a rigorous and collaborative way.

\section*{2. LLM for KG Reasoning}
My future research goal is to build LLM systems that reason more reliably over large knowledge graphs, especially in biomedicine~\cite{llmkgsurvey,kgagentmed}. Many biomedical questions require multi-hop composition across genes, proteins, pathways, phenotypes, drugs, diseases, and clinical concepts. In principle, however, the space of possible reasoning chains grows combinatorially: for a graph with average degree $d$ and hop depth $k$, naive path exploration yields on the order of $d^k$ candidates. Existing systems therefore rely on shallow retrieval, heuristic pruning, or manually chosen seed nodes. These approximations are often useful, but they become brittle exactly where scientific discovery demands robustness.

At the same time, current LLMs are not yet strong graph reasoners on their own. Even frontier models continue to struggle with faithful traversal and composition on graph-style tasks; Anthropic's April 16, 2026 Claude Opus 4.7 report still shows limited performance on GraphWalks search and parent-retrieval tasks~\cite{anthropicopus}. This matches what I observed in \textit{GraphText}~\cite{graphtext}: LLMs are often good at summarizing local evidence, but much less reliable when they must systematically search a structured space. For this reason, I do not think the right objective is simply to ``let the LLM reason over the KG.'' The key challenge is to design a better division of labor between graph-native and language-native reasoning modules.

My central hypothesis is that LLMs should reason \textit{over explicit path evidence}, rather than search a knowledge graph from scratch. In this view, a graph model first proposes a small set of reliable \textit{path instances} or reasoning traces, and the LLM then analyzes, compares, aggregates, and verbalizes them. This is a different interface from node-only retrieval or simply stuffing a subgraph into context. Paths are attractive because they are structured, compositional, and interpretable: they expose how evidence flows across the graph and which relations support an answer. In biomedicine, this matters because an answer without a mechanism is often much less useful than one tied to a plausible chain of evidence.

I would like to pursue this agenda through three connected research thrusts. The first is \textit{tractable path proposal over holistic knowledge graphs}. Here, I want to study how graph models --- including GNNs, knowledge-graph embedding models, and logical query encoders~\cite{logicalqueries} --- can propose informative multi-hop paths without enumerating the full search space. A key idea is that graph models can be fine-tuned over the entire knowledge base in a way that current LLMs cannot, allowing them to encode global graph regularities, relation semantics, and structural shortcuts. I am especially interested in adaptive breadth-depth retrieval strategies, inspired by recent work such as ARK from Zitnik Lab~\cite{ark}, but extending them from node exploration to path-level proposal and ranking.

The second thrust is \textit{faithful language reasoning over candidate paths}. Once a graph module proposes a small set of path instances, an LLM can be used for semantic comparison, evidence aggregation, contradiction detection, uncertainty expression, and natural-language explanation. I want to develop objectives and evaluation protocols for \textit{path faithfulness}: a model should not only produce a correct answer, but should also ground that answer in explicit graph evidence that survives perturbation and scrutiny. This may involve contrastive training over plausible and implausible paths, verifier models that check relational consistency, or distillation pipelines in which stronger reasoners teach smaller models to select better evidence.

The third thrust is \textit{biomedical and multimodal KG reasoning}. In realistic scientific settings, knowledge graphs do not exist in isolation; they interact with literature, omics measurements, clinical time series, ontologies, and experimental observations. I therefore want to study how path-grounded KG reasoning can be integrated with multimodal evidence, so that the system can move between structured biomedical knowledge and unstructured scientific context. Example applications include mechanism-aware biomedical question answering, candidate target prioritization, therapy repurposing, and hypothesis generation where the output is not just a ranked entity but a small set of evidence-backed mechanistic routes.

I am particularly excited about pursuing this work in Zitnik Lab because the lab uniquely combines large-scale biomedical knowledge graphs, geometric learning, multimodal foundation models, and agentic retrieval systems with a strong orientation toward scientific impact. My hope is to contribute graph-centered methodology while learning how to push it toward more consequential biomedical problems. In the long run, I want to help build knowledge-grounded foundation models that can move fluidly between graph structure, language, and scientific evidence --- models that support discovery not only by making predictions, but by exposing the reasoning paths that make those predictions actionable.

\clearpage
\renewcommand{\refname}{3. References}
\scriptsize
\begin{thebibliography}{99}\setlength{\itemsep}{0.16em}
\bibitem{genezip} J. Zhao, X. Liu, Z. Zhan, X. Yuan, H. Guo, and J. Tang. \textit{GeneZip: Region-Aware Compression for Long Context DNA Modeling}. arXiv 2026.
\bibitem{raftppi} J. Zhao, Z. Zhan, N. Chaudhary, X. Yuan, Z. Zhang, Q. Cong, J. Zhou, S. Misra, and J. Tang. \textit{Fast Proteome-Scale Protein Interaction Retrieval via Residue-Level Factorization}. ICLR 2026.
\bibitem{hax} H. Mao, Z. Chen, W. Tang, J. Zhao, Y. Ma, T. Zhao, N. Shah, M. Galkin, and J. Tang. \textit{Overcoming Long-Context Limitations of State-Space Models via Context-Dependent Sparse Attention}. NeurIPS 2025.
\bibitem{pandemicllm} H. Du, Y. Zhao, J. Zhao, S. Xu, X. Lin, Y. Chen, L. Gardner, and H. F. Yang. \textit{Advancing real-time infectious disease forecasting using large language models}. Nature Computational Science 2025.
\bibitem{graphany} J. Zhao, Z. Zhu, M. Galkin, H. Mostafa, M. Bronstein, and J. Tang. \textit{Fully-inductive Node Classification on Arbitrary Graphs}. ICLR 2025.
\bibitem{graphtext} J. Zhao, L. Zhuo, Y. Shen, M. Qu, K. Liu, M. Bronstein, Z. Zhu, and J. Tang. \textit{GraphText: Graph Reasoning in Text Space}. arXiv 2023; NeurIPS AFM Workshop 2024.
\bibitem{glem} J. Zhao, M. Qu, C. Li, H. Yan, Q. Liu, R. Li, X. Xie, and J. Tang. \textit{Learning on Large-scale Text-attributed Graphs via Variational Inference}. ICLR 2023.
\bibitem{gsr} J. Zhao, Q. Wen, M. Ju, C. Zhang, and Y. Ye. \textit{Self-Supervised Graph Structure Refinement for Graph Neural Networks}. WSDM 2023.
\bibitem{house} R. Li, J. Zhao, C. Li, D. He, Y. Wang, Y. Liu, H. Sun, S. Wang, W. Deng, Y. Shen, X. Xie, and Q. Zhang. \textit{HousE: Knowledge Graph Embedding with Householder Parameterization}. ICML 2022.
\bibitem{gophormer} J. Zhao, R. Li, Q. Wen, C. Li, Y. Wang, Y. Liu, H. Sun, X. Xie, and Y. Ye. \textit{Gophormer: Ego-Graph Transformer for Node Classification}. arXiv 2021.
\bibitem{mvse} J. Zhao, Q. Wen, S. Sun, Y. Ye, and C. Zhang. \textit{Multi-View Self-Supervised Heterogeneous Graph Embedding}. ECML/PKDD 2021.
\bibitem{hgsl} J. Zhao, X. Wang, C. Shi, B. Hu, G. Song, and Y. Ye. \textit{Heterogeneous Graph Structure Learning for Graph Neural Networks}. AAAI 2021.
\bibitem{nshe} J. Zhao, X. Wang, C. Shi, Z. Liu, and Y. Ye. \textit{Network Schema Preserving Heterogeneous Information Network Embedding}. IJCAI 2020.
\bibitem{logicalqueries} W. Hamilton, P. Bajaj, M. Zitnik, D. Jurafsky, and J. Leskovec. \textit{Embedding Logical Queries on Knowledge Graphs}. NeurIPS 2018.
\bibitem{llmkgsurvey} J. Pan, S. Razniewski, J.-C. Kalo, S. Singhania, J. Chen, S. Dietze, and others. \textit{Large Language Models and Knowledge Graphs: Opportunities and Challenges}. arXiv 2023.
\bibitem{kgagentmed} X. Su, Y. Wang, S. Gao, X. Liu, V. Giunchiglia, D.-A. Clevert, and M. Zitnik. \textit{Knowledge Graph Based Agent for Complex, Knowledge-Intensive QA in Medicine}. arXiv 2024.
\bibitem{ark} J. Polonuer, L. Vittor, I. Arango, A. Noori, D. A. Clifton, L. Del Corro, and M. Zitnik. \textit{Autonomous Knowledge Graph Exploration with Adaptive Breadth-Depth Retrieval}. arXiv 2026.
\bibitem{anthropicopus} Anthropic. \textit{Introducing Claude Opus 4.7}. Product announcement, April 16, 2026.
\end{thebibliography}

\end{document}
